\documentclass{article}
\usepackage{amsmath,amssymb}
\usepackage{booktabs}
\usepackage{graphicx}
\usepackage{hyperref}
\usepackage{natbib}

% sets
\newcommand{\mA}{\mathcal{A}}
\newcommand{\mD}{\mathcal{D}}
\newcommand{\mE}{\mathcal{E}}
\newcommand{\mN}{\mathcal{N}}
\newcommand{\mP}{\mathcal{P}}
\newcommand{\mV}{\mathcal{V}}

% vectors
\newcommand{\bw}{\boldsymbol{w}}
\newcommand{\bx}{\boldsymbol{x}}
\newcommand{\bz}{\boldsymbol{z}}


\title{Experimental Roadmap for Validating Fenchel--Young Surrogates in Kidney Exchange}
\author{Weikang}
\date{}

\begin{document}
\maketitle

\section{Motivation}

The purpose of the experiments is not only to compare decision-focused learning
(DFL) against a two-stage predict-then-optimize pipeline, but also to establish
whether the Fenchel--Young (FY) perturbed-optimizer objective
\citep{berthet2020learning} is a meaningful surrogate for the target decision
loss in our kidney exchange setting.

This distinction matters because the realistic synthetic benchmark used in our
main experiments has a complex reward mechanism. The ground-truth edge reward is
not a simple linear function of utility, but a nonlinear product of transplant
success probability, QALY, scarcity priority, graph-structural variables, and
deterministic edge-level noise. Therefore, a direct comparison between DFL and a
two-stage baseline on this benchmark alone cannot isolate the behavior of the
FY objective. If DFL improves the final matching score, the improvement may be
due to the surrogate, the choice of perturbation scale, the warm start, or model
capacity. If DFL does not improve, the failure may instead be caused by model
misspecification or an overly complex label-generation process.

We therefore separate the experimental study into several layers, moving from a
controlled noisy-linear synthetic reward benchmark to the full realistic
synthetic benchmark.

\clearpage
\begin{figure}[p]
    \centering
    \includegraphics[width=0.90\textwidth]{surrogate_experiment_results/Step1/epsilon=0.2/trajectory_contour.png}\\[-0.45em]
    \includegraphics[width=0.90\textwidth]{surrogate_experiment_results/Step1/epsilon=0.2/trajectory_3d_metrics.png}\\[-0.45em]
    \includegraphics[width=0.90\textwidth]{surrogate_experiment_results/Step1/epsilon=1.0/trajectory_contour.png}\\[-0.45em]
    \includegraphics[width=0.90\textwidth]{surrogate_experiment_results/Step1/epsilon=1.0/trajectory_3d_metrics.png}\\[-0.35em]
    \caption{Step~1 diagnostics on a controlled noisy-linear synthetic reward
    benchmark with \(\sigma=0.10\). The four panels, from top to bottom, show
    the \(\epsilon=0.2\) synthetic-label decision-gap landscape,
    \(\epsilon=0.2\) trajectory metrics, the \(\epsilon=1.0\)
    synthetic-label decision-gap landscape, and \(\epsilon=1.0\) trajectory
    metrics. The decision gap is the oracle objective gap under the synthetic
    label weights, not a clinical regret measure. All panels use the same 100
    graph instances from
    \texttt{dataset/processed/step1\_noisy\_linear\_sigma010}.}
    \label{fig:step1-diagnostics-all}
\end{figure}
\clearpage

\section{Experimental Roadmap}

\subsection{Step 1: Surrogate Validity Under a Noisy Linear Synthetic Reward}

The first experiment should validate whether the perturbed FY objective is
directionally aligned with downstream decision quality in a controlled setting.
The current Step~1 run uses a noisy linear synthetic reward:
\[
    w^{\mathrm{syn}}_e
    =
    \max\left\{
        0,\,
        (10u_e + 5c_e)(1+\xi_e)
    \right\},
\]
where \(u_e\) is the edge utility and \(c_e\) is the recipient cPRA associated
with edge \(e\). The deterministic edge-level noise \(\xi_e\) is generated by
the processed batch
\texttt{dataset/processed/step1\_noisy\_linear\_sigma010}, with multiplicative
noise scale \(\sigma=0.10\). The corresponding prediction model remains the
two-feature linear probe
\[
    \hat w_e(\theta) = \theta_1 u_e + \theta_2 c_e.
\]

This setting preserves a transparent two-feature signal while adding controlled
edge-level variation that the two-feature probe cannot perfectly fit. We
therefore evaluate parameters by the \emph{synthetic-label decision gap}:
\[
    G(\theta)
    =
    \frac{1}{N}
    \sum_{i=1}^{N}
    \left[
        (w_i^{\mathrm{syn}})^\top y_i^\star
        -
        (w_i^{\mathrm{syn}})^\top S(\hat w_i(\theta))
    \right],
\]
where \(y_i^\star=S(w_i^{\mathrm{syn}})\) is the KEP solution optimized under
the synthetic label weights. This is an oracle objective gap under the synthetic
reward labels; it should not be read as regret with respect to real transplant
outcomes, clinical utility, or patient welfare. If the FY surrogate is
appropriate, its trajectory should still move toward low decision-gap regions
under the synthetic reward labels, even though exact reward recovery is no
longer possible.

The main diagnostic objects are:
\begin{itemize}
    \item the synthetic-label decision-gap landscape \(G(\theta_1,\theta_2)\);
    \item the perturbed FY objective landscape \(L_{\mathrm{FY}}(\theta_1,\theta_2)\);
    \item the two-stage MSE solution \(\theta^\star_{\mathrm{2s}}\);
    \item the synthetic-label oracle solution \(y^\star=S(w^{\mathrm{syn}})\);
    \item the best probe parameter on the grid
    \[
        \theta_{\mathrm{gap\text{-}min}}
        =
        \arg\min_{\theta \in \mathrm{grid}} G(\theta);
    \]
    \item the FY minimizer
    \[
        \theta^\star_{\mathrm{FY}}(\epsilon)
        =
        \arg\min_{\theta} L_{\mathrm{FY},\epsilon}(\theta).
    \]
\end{itemize}

The primary metric should be decision gap rather than parameter distance:
\[
    \Delta G(\epsilon)
    =
    G(\theta^\star_{\mathrm{FY}}(\epsilon))
    -
    \min_{\theta} G(\theta).
\]
This is important because the KEP solver is piecewise constant, so the
decision-gap landscape often has a low-gap plateau rather than a unique isolated
minimizer. In such cases, parameter-space distance to one selected grid point is
not well-defined, whereas \(\Delta G\) remains meaningful.

This step answers the question:

\begin{quote}
Does the FY surrogate point toward low synthetic-label decision-gap regions when
the predictive model captures the main noise-free linear signal but cannot represent
deterministic edge-level noise?
\end{quote}

In the current Step~1 diagnostic, this question is tested on 100 graph
instances, using the same sampled instances for the MSE and FY trajectories. We
track the full optimization trajectories of both the two-stage MSE objective and
the decision-focused FY objective in the \((\theta_1,\theta_2)\) plane, and
evaluate each point by synthetic-label decision gap. The saved trajectories
contain \(1001\) points
\((1000\) training epochs plus the initialization\()\). The initial parameter
has average decision gap \(3.5402\).

The two-stage MSE trajectory still recovers the noise-free linear component of the noisy
label, converging to approximately \((10.00,5.01)\). However, because of the
deterministic multiplicative noise, the final MSE solution does not reach zero
decision gap; its final average decision gap is \(0.8759\). The FY trajectories
show a different, decision-focused behavior. They do not need to recover the
noise-free linear coefficient \((10,5)\). Instead, they move through regions
that improve the induced KEP decisions under the synthetic label weights. For
\(\epsilon=0.2\), the trajectory reaches its best observed decision gap early,
at epoch \(42\), before continuing to reduce the perturbed FY objective. For
\(\epsilon=1.0\), the trajectory grows to a larger parameter scale and reaches
its best observed decision gap later, around epoch \(781\).

\begin{center}
\small
\resizebox{\textwidth}{!}{%
\begin{tabular}{@{}cccccc@{}}
\toprule
method & \(\epsilon\) & final \(\theta\) & final FY objective & final decision gap & best decision gap \\
\midrule
MSE & -- & \((10.0017,\,5.0147)\) & -- & \(0.8759\) & \(0.8759\) \\
FY & \(0.2\) & \((2.8135,\,1.3589)\) & \(0.3609\) & \(0.8714\) & \(0.7183\) (ep. \(42\)) \\
FY & \(1.0\) & \((13.9796,\,6.7825)\) & \(1.8042\) & \(0.8714\) & \(0.8668\) (ep. \(781\)) \\
\bottomrule
\end{tabular}
}
\end{center}

These results support the intended interpretation of the diagnostic. The
two-stage method identifies the noise-free linear signal in the noisy label, while
the FY surrogate searches in a decision-induced geometry. In this noisy
controlled setting, the FY trajectory can visit lower-gap regions than the
MSE endpoint, but the final FY endpoint is not necessarily the best-gap
iterate. This distinction is important for interpreting FY optimization: the
perturbed FY objective may continue to optimize perturbation robustness even
after the deterministic KEP decision has already reached a low-gap region.

This leads to the following practical observations.
\begin{enumerate}
    \item \textbf{The FY objective is decision-oriented rather than reward-fitting.}
    Unlike the two-stage MSE trajectory, which recovers the noise-free linear component
    \((10,5)\), the FY trajectory moves according to changes in the induced KEP
    decision. With \(\epsilon=0.2\), the best observed FY point has decision gap
    \(0.7183\), below the final MSE decision gap \(0.8759\), despite being far
    from \((10,5)\) in parameter space.

    \item \textbf{The perturbation scale controls decision stability and
    parameter scale.}
    With \(\epsilon=0.2\), the final FY endpoint is
    \((2.8135,1.3589)\). With \(\epsilon=1.0\), the final endpoint is
    \((13.9796,6.7825)\). The larger perturbation scale encourages a larger
    prediction scale, because the same decision must remain stable under
    stronger perturbations. This is consistent with the FY gradient form
\[
    \nabla_{\hat w} L_{\mathrm{FY}}(\hat w; y^\star)
    =
    \mathbb{E}_{z}\!\left[y^\star(\hat w + z)\right] - y^\star,
\]
    which penalizes deviations between the perturbation-averaged optimizer
    output and the synthetic-label oracle decision.

    \item \textbf{The best-gap point need not be the final FY endpoint.}
    For \(\epsilon=0.2\), the minimum observed decision gap is \(0.7183\) at
    epoch \(42\), while the final decision gap is \(0.8714\). For
    \(\epsilon=1.0\), the best observed decision gap is \(0.8668\) at epoch
    \(781\), and the final decision gap is \(0.8714\). This shows that in the
    noisy setting, early stopping or model selection by decision gap can matter.

    \item \textbf{Zero decision gap is unrealistic for this two-feature probe.}
    The MSE endpoint is close to \((10,5)\), but its final decision gap remains
    \(0.8759\). Thus the current Step~1 is no longer a pure sanity check for
    recovering the noise-free coefficient; it is a controlled
    decision-quality diagnostic under mild label noise.
\end{enumerate}

\subsection{Step 2: Controlled Model Misspecification}

The second experiment should introduce a limited amount of misspecification.
For example, the synthetic label reward may include an additional feature or
interaction:
\[
    w^{\mathrm{syn}}_e
    =
    a u_e + b c_e + d\,\mathrm{age}^{\mathrm{donor}}_e,
\]
or
\[
    w^{\mathrm{syn}}_e
    =
    a u_e + b c_e + d\,u_e c_e.
\]
The prediction model can then be kept deliberately simpler:
\[
    \hat w_e(\theta) = \theta_1 u_e + \theta_2 c_e.
\]

This setting tests whether FY optimization is robust when the model cannot
perfectly recover the edge rewards. The question is no longer whether FY can
recover the synthetic label function, but whether it can still produce better
decisions than an MSE-trained predictor.

The same landscape diagnostics should be reported:
\[
    G(\theta), \quad
    L_{\mathrm{FY},\epsilon}(\theta), \quad
    G(\theta^\star_{\mathrm{2s}}), \quad
    G(\theta^\star_{\mathrm{FY}}), \quad
    \Delta G.
\]

This step answers the question:

\begin{quote}
When the predictive model is misspecified, does the FY surrogate still move the
parameters toward lower synthetic-label decision gap than the two-stage MSE
solution?
\end{quote}

\subsection{Step 3: Training-Level DFL vs Two-Stage Comparison}

After the landscape-level surrogate diagnostics, we should run actual training
experiments on the same controlled settings. Both methods should use the same
train/validation/test split, the same model class, and the same KEP solver.

The two-stage baseline first minimizes prediction error:
\[
    \min_{\theta}
    \sum_{e}
    \left(\hat w_e(\theta)-w^{\mathrm{syn}}_e\right)^2,
\]
then solves the KEP using the predicted weights \(\hat \bw(\theta)\).

The DFL model starts either from the two-stage checkpoint or from random
initialization and then updates the predictor using the FY gradient estimate:
\[
    \nabla_{\hat \bw} L_{\mathrm{FY}}
    \approx
    \bar y_{\epsilon}(\hat \bw) - y^\star,
\]
where
\[
    \bar y_{\epsilon}(\hat \bw)
    =
    \frac{1}{M}
    \sum_{m=1}^{M}
    S(\hat \bw+\epsilon z^{(m)}),
    \qquad
    z^{(m)}\sim \mathcal{N}(0,I).
\]

The primary evaluation metric should be final achieved synthetic-reward score,
or equivalently average synthetic-label decision gap:
\[
    G(\theta)
    =
    \frac{1}{N}
    \sum_{i=1}^{N}
    \left[
        (\bw_{\mathrm{syn}}^i)^\top y^\star(\bw_{\mathrm{syn}}^i)
        -
        (\bw_{\mathrm{syn}}^i)^\top S(\hat \bw^i(\theta))
    \right].
\]

This step answers the question:

\begin{quote}
Does the landscape-level alignment between FY and synthetic-label decision gap
translate into a training-level improvement over the two-stage baseline?
\end{quote}

\subsection{Step 4: Realistic Synthetic Benchmark}

The final main benchmark should use the currently processed dataset, now named
\texttt{dataset/processed/realistic\_synthetic\_benchmark}. This benchmark uses
the nonlinear synthetic medical reward:
\[
    w^{\mathrm{syn}}_e
    =
    E[\mathrm{transplant}]_e
    \cdot
    \mathrm{QALY}_e
    \cdot
    \mathrm{priority}_e
    \cdot
    (1+\xi_e).
\]
This label incorporates utility, recipient cPRA, donor and recipient age, blood
type context, arc failure probability, vertex failure probability, scarcity,
donor flexibility, reciprocal edges, and deterministic edge-level noise.

This benchmark should be described as a realistic synthetic benchmark rather
than a simple surrogate-identification experiment. Its role is to test whether
the DFL pipeline improves exchange-plan quality in a richer synthetic
environment.

The interpretation should be careful:
\begin{itemize}
    \item If DFL improves over two-stage, this provides empirical evidence that
    the FY-based pipeline can help in a complex KEP setting.
    \item If DFL does not improve, this does not by itself prove that the FY
    surrogate is invalid; the failure may come from model misspecification,
    insufficient features, perturbation scale, or optimization instability.
\end{itemize}

This step answers the question:

\begin{quote}
Under a medically motivated synthetic reward model, does FY-based DFL improve
the final kidney exchange objective compared with the two-stage pipeline?
\end{quote}

\subsection{Step 5: Sensitivity and Ablation Studies}

The last layer should quantify which design choices affect the result. The most
important ablations are:
\begin{itemize}
    \item perturbation scale \(\epsilon\);
    \item number of perturbation samples \(M\);
    \item warm start from the two-stage solution versus random initialization;
    \item feature set: utility only, utility+cPRA, small linear features,
    failure context features, and full features;
    \item label noise on versus off;
    \item solver formulation, such as CF, PIEF, and hybrid.
\end{itemize}

For the perturbation scale, it is important to distinguish absolute and relative
noise. Since the KEP solver is an argmax, its output is invariant to positive
rescaling of all weights. Therefore the effective smoothing strength is governed
by the relative scale of the perturbation:
\[
    S(\hat \bw+\epsilon z)
    =
    S\left(\frac{\hat \bw}{\epsilon}+z\right).
\]
Thus, \(\epsilon\) should be interpreted relative to the scale of
\(\hat \bw\). In the landscape plots, \(\epsilon\) may be calibrated relative to
the MSE weight scale, while in FY training it may be used as an absolute
noise standard deviation. This difference should be reported explicitly.

\section{Recommended Presentation Order}

The final paper should present the experiments in the following order:
\begin{enumerate}
    \item noisy-linear synthetic reward surrogate-validity benchmark;
    \item controlled misspecification benchmark;
    \item training-level DFL versus two-stage comparison on controlled labels;
    \item realistic synthetic benchmark;
    \item epsilon, sample-size, feature-set, and initialization ablations.
\end{enumerate}

This ordering separates mechanism validation from final empirical performance.
It allows the paper to make precise claims: first, whether the FY surrogate is
aligned with synthetic-label decision gap under controlled conditions; second, whether this
alignment survives moderate misspecification; and finally, whether it yields
practical improvements on the realistic synthetic KEP benchmark.

\bibliographystyle{plainnat}
\bibliography{ref}

\end{document}
